\documentclass[10pt,letterpaper]{article}
\usepackage[T1]{fontenc}
\usepackage[utf8]{inputenc}
\usepackage[sfdefault]{sourcesanspro}
\usepackage{inconsolata}
\usepackage[margin=0.84in,headheight=15pt]{geometry}
\usepackage{microtype}
\usepackage{xcolor}
\usepackage{booktabs,longtable,tabularx,array}
\usepackage{enumitem}
\usepackage{fancyhdr,lastpage}
\usepackage[most]{tcolorbox}
\usepackage{listings}
\usepackage{titlesec}
\usepackage{hyperref}
\usepackage{amsmath,amssymb}
\usepackage{multicol}
\definecolor{navy}{HTML}{15324C}
\definecolor{teal}{HTML}{087E83}
\definecolor{slate}{HTML}{536879}
\definecolor{paper}{HTML}{F2F6F8}
\definecolor{warm}{HTML}{E8F4F3}
\hypersetup{colorlinks=true,urlcolor=teal,linkcolor=navy,pdfauthor={Software Engineering Reference Suite},pdfsubject={MISRA C++:2023 engineering guidance}}
\titleformat{\section}{\Large\bfseries\color{navy}}{\thesection}{0.7em}{}
\titleformat{\subsection}{\large\bfseries\color{teal}}{\thesubsection}{0.6em}{}
\titleformat{\subsubsection}{\normalsize\bfseries\color{navy}}{\thesubsubsection}{0.5em}{}
\setlist[itemize]{leftmargin=*,topsep=2pt,itemsep=2pt}
\setlist[enumerate]{leftmargin=*,topsep=2pt,itemsep=3pt}
\setlength{\parskip}{5pt}
\setlength{\parindent}{0pt}
\setlength{\emergencystretch}{2em}
\pagestyle{fancy}
\fancyhf{}
\fancyhead[L]{\small\color{slate}MISRA C++:2023 Engineering Suite}
\fancyhead[R]{\small\color{slate}\nouppercase{\leftmark}}
\fancyfoot[L]{\scriptsize\color{slate}Independent engineering guidance; see licensed source}
\fancyfoot[R]{\scriptsize\color{slate}\thepage\ / \pageref{LastPage}}
\newtcolorbox{sourcebox}{colback=paper,colframe=navy!30,boxrule=0.45pt,arc=1.6pt,breakable}
\newtcolorbox{practicebox}{colback=warm,colframe=teal!40,boxrule=0.45pt,arc=1.6pt,breakable}
\newcommand{\SourceNote}[1]{\begin{sourcebox}\textbf{SOURCE BASIS / NORMATIVE POINTER}\quad #1\end{sourcebox}}
\newcommand{\PracticeNote}[1]{\begin{practicebox}\textbf{ENGINEERING RECOMMENDATION / NOT MISRA TEXT}\quad #1\end{practicebox}}
\newcommand{\Record}[1]{\textsf{\bfseries #1}}
\lstdefinestyle{cppsuite}{language=C++,basicstyle=\ttfamily\footnotesize,keywordstyle=\color{navy}\bfseries,commentstyle=\color{slate},stringstyle=\color{teal},breaklines=true,breakatwhitespace=true,columns=fullflexible,frame=single,framerule=0.3pt,rulecolor=\color{navy!35},backgroundcolor=\color{paper},showstringspaces=false,tabsize=2,xleftmargin=0.5em,xrightmargin=0.5em,aboveskip=4pt,belowskip=4pt}
\lstset{style=cppsuite}
\newcommand{\SuiteTitle}[3]{%
\begin{titlepage}\thispagestyle{empty}
\vspace*{1.5cm}
{\color{teal}\rule{2.4cm}{2pt}}\par\vspace{1cm}
{\Huge\bfseries\color{navy}MISRA C++:2023\\[0.4em]#1}\par\vspace{0.7cm}
{\Large\color{slate}#2}\par\vspace{1.2cm}
\begin{sourcebox}\textbf{Source baseline:} MISRA C++:2023, Guidelines for the use of C++17 in critical systems, October 2023 (user-supplied licensed copy). This is an independent companion, not an official MISRA publication, rule replacement, certified checker, or claim of project compliance.\end{sourcebox}
\vfill{\large\color{navy}Technical edition \textbullet\ 16 October 2026}\par\smallskip
{\small\color{slate}#3}\par\end{titlepage}
\tableofcontents\clearpage
}
\newcommand{\StandardsBoundary}{%
\section*{How to read normative status}\addcontentsline{toc}{section}{How to read normative status}
\textbf{Normative source:} the licensed MISRA C++:2023 guideline text, including any applicable \emph{Amplification} and \emph{Exception}. A guideline \emph{headline} alone is not a complete compliance test. The source's Section~3 defines classifications, categories, decidability, analysis scope, and applicability.\par
\textbf{Independent guidance:} all review questions, programming patterns, lifecycle processes, sample thresholds, ownership assignments, controls, and CI configurations in this suite are \emph{engineering recommendations}, not statements issued by MISRA.\par
\textbf{Critical limitation:} no individual example is asserted compliant with the whole standard. For formal compliance claims, follow the applicable edition of \emph{MISRA Compliance} as the source itself instructs, with a documented compliance plan and supporting evidence.\par
}

\begin{document}
\SuiteTitle{Deviation and Disposition Management}{Formal Required deviations, Advisory records, exception decisions}{Volume 08 / Risk Governance}
\StandardsBoundary

\section{Deviation semantics}
A deviation is not a static-analyzer suppression, an annotation in source code, or a reviewer belief that the risk is small. In the source model, an actual violation of a Required guideline must have a formal, justified deviation to support a claim. Mandatory guidelines cannot be deviated. Advisory violations require documentation even when the formal deviation process is not used.

A source-defined Exception is different: the specified case is outside the guideline applicability and does not create a violation. A false-positive diagnostic is also not an actual violation once the supporting semantic argument is substantiated.


\section{DVR-01 decision tree}
\begin{tcolorbox}[colback=paper,colframe=navy!25,title={Reviewer sequence}]
1. Verify exact guideline ID, edition, target configuration and source location.\\
2. Does the official amplification or source-defined exception change applicability? If so, document the proof and close as no violation.\\
3. Is the diagnostic demonstrably false? Document a reproducible tool finding and corrective action to tool mapping.\\
4. Is an actual violation present?\\
\hspace*{1em} Mandatory: repair; no deviation.\\
\hspace*{1em} Required: repair or request formal deviation.\\
\hspace*{1em} Advisory: repair or provide documented justification/disapplication (where allowed).\\
5. Independently validate acceptance, residual risk and release scope.
\end{tcolorbox}

\section{Proposed deviation request fields}
\small\setlength{\tabcolsep}{3pt}\renewcommand{\arraystretch}{1.15}
\begin{longtable}{p{0.22\linewidth}p{0.66\linewidth}}
\toprule
Field & Minimum content\\ \midrule
\endfirsthead
\toprule
Field & Minimum content\\ \midrule
\endhead
Identifier & DVR ID; guideline ID and edition; source file, line, revision and component.\\
Category & Required; verify not Mandatory. Stronger project recategorization is considered.\\
Trigger & Finding reference, build configuration, analyzer and tool version.\\
Exact nonconformance & Specific behavior and applicable amplification/exception analysis.\\
Necessity & Technical or system requirement that prevents straightforward correction.\\
Alternatives & Feasible fixes, why rejected, impact of each alternative.\\
Hazard and risk & Failure mechanism, severity, likelihood reasoning, affected interfaces and operation.\\
Compensating verification & Tests, runtime checks, design constraints, additional static/dynamic evidence.\\
Bounded scope & Permitted sites/configurations/releases and external dependency assumptions.\\
Approval & Independent engineer and designated safety/security authority, dates and signoff.\\
Lifecycle & Review date, expiration/trigger, ownership and closure evidence.\\
\bottomrule
\end{longtable}
\normalsize

\section{Required vs advisory records}
For a Required actual violation, use a formal DVR that documents technical necessity and the controlled exception decision. For Advisory actual violations, use an advisory assessment record, unless the project chooses to run the same formal process. Source-defined advisory disapplication still needs documented project justification.

Avoid business-process statuses such as ``Won't Fix as the entire explanation; tie every outcome to a guideline category and governing decision.


\section{Worked scenario: constrained C++17 target}
Assume a third-party component uses a C++ facility restricted by a Required MISRA rule and refactoring would break a verified interface. The engineer first demonstrates the actual applicability and scope of the violation; asks whether removing the code from the shipped binary or using a source exception eliminates the issue; then documents safe alternatives and any residual risk. The approval is limited to the exact component and target, not permission to violate the rule anywhere.

This hypothetical scenario does not establish that any particular library usage merits approval. Source conditions, target-specific behavior and the applicable MISRA Compliance process govern the actual outcome.


\section{Metrics and expiry controls}
\small\setlength{\tabcolsep}{3pt}\renewcommand{\arraystretch}{1.15}
\begin{longtable}{p{0.28\linewidth}p{0.6\linewidth}}
\toprule
Metric & Recommended interpretation\\ \midrule
\endfirsthead
\toprule
Metric & Recommended interpretation\\ \midrule
\endhead
Open Required deviations & Count individually and by release; age does not make them acceptable.\\
Mandatory violations & Always list separately; a deviation count cannot offset one.\\
Advisory justifications & Maintain evidence and project-recategorization reason.\\
Expired deviations & Treat as unauthorized until renewed or fixed.\\
False positives & Track tool issues separately from actual violations and formal deviations.\\
Coverage gaps & Communicate unverified rules and variants rather than manipulating pass rates.\\
\bottomrule
\end{longtable}
\normalsize

\end{document}
